\documentclass[uplatex,dvipdfmx,a4paper,10pt]{jsarticle}

\usepackage[margin=20mm,headheight=15pt,headsep=7mm,footskip=12mm]{geometry}
\usepackage[dvipdfmx]{graphicx}
\usepackage[dvipdfmx]{xcolor}
\usepackage[haranoaji]{pxchfon}
\definecolor{mainblue}{HTML}{1F4E79}
\definecolor{lightblue}{HTML}{EAF3F8}
\definecolor{lightgray}{HTML}{F5F5F5}
\definecolor{darkgray}{HTML}{444444}
\usepackage[dvipdfmx]{hyperref}
\usepackage{pxjahyper}
\hypersetup{colorlinks=true,linkcolor=mainblue,urlcolor=mainblue,citecolor=mainblue}
\usepackage{fancyhdr}
\usepackage{longtable,booktabs,tabularx,array}
\usepackage{enumitem}
\usepackage[most]{tcolorbox}
\usepackage{amsmath}
\usepackage{tikz}
\usetikzlibrary{arrows.meta,positioning,fit,shapes.geometric}

\setlength{\parindent}{0pt}
\setlength{\parskip}{0.55em}
\renewcommand{\arraystretch}{1.25}
\newcolumntype{Y}{>{\raggedright\arraybackslash}X}
\newcolumntype{L}[1]{>{\raggedright\arraybackslash}p{#1}}
\sloppy
\emergencystretch=3em

\pagestyle{fancy}
\fancyhf{}
\fancyhead[L]{\small 06 章 ソフトウェアエンジニアリング運用 SWEBOK v4.0a 日本語まとめ}
\fancyhead[R]{\small \thepage}
\renewcommand{\headrulewidth}{0.4pt}
\fancyfoot[C]{\small \thepage/\pageref{LastPage}}
\usepackage{lastpage}

\newtcolorbox{pointbox}[1]{
  breakable,
  colback=lightblue,
  colframe=mainblue,
  fonttitle=\bfseries,
  title={#1},
  left=1.5mm,right=1.5mm,top=1mm,bottom=1mm,
  boxrule=0.5pt
}
\newtcolorbox{notebox}[1]{
  breakable,
  colback=lightgray,
  colframe=darkgray,
  fonttitle=\bfseries,
  title={#1},
  left=1.5mm,right=1.5mm,top=1mm,bottom=1mm,
  boxrule=0.4pt
}
\newcommand{\todoprompt}[2]{%
\begin{notebox}{TODO（図）：#1}
#2

\textbf{画像生成用プロンプト}

白背景、モノクロ、A4本文幅に収まる学習用図解を作成する。#2 日本語ラベル、細線、講義ノート風、余白を広め、印刷で読みやすい図にする。
\end{notebox}
}

\begin{document}

\begin{titlepage}
  \vspace*{10mm}
  {\sffamily\bfseries\fontsize{24}{28}\selectfont 06 章\par}
  \vspace{3mm}
  {\sffamily\bfseries\fontsize{30}{34}\selectfont ソフトウェアエンジニアリング運用\par}
  \vspace{2mm}
  {\Large Software Engineering Operations の日本語まとめ\par}
  \vspace{3mm}
  {\large SWEBOK Guide v4.0a Chapter 06 を初学者向けに再構成\par}
  \vspace{4mm}
  作成日：2026年4月29日

  \vspace{8mm}
  \begin{pointbox}{この資料の主張}
  ソフトウェアエンジニアリング運用は、完成したソフトウェアを「本番に置く作業」だけではない。運用は、配備、設定、監視、障害対応、変更管理、容量管理、バックアップ、災害復旧、サービス水準、セキュリティ、自動化を含む一連の工学活動である。現代の運用では、開発と運用を分離せず、開発運用連携（DevOps）、インフラストラクチャのコード化（Infrastructure as Code, IaC）、サイト信頼性エンジニアリング（Site Reliability Engineering, SRE）などを使い、ソフトウェアを安全かつ継続的に改善することが重要である。
  \end{pointbox}

  \begin{pointbox}{この資料の読み方}
  専門用語は原則として日本語で示し、必要に応じて英語を括弧内に併記する。英語名は、元資料、標準、ツール名を調べる手がかりとして残す。初学者が前提知識なしに読めるように、各節では「何のための概念か」「実務では何を見るべきか」「失敗すると何が起こるか」を重視する。文体は、だである調で統一する。
  \end{pointbox}

  \vfill
  \begin{center}
  \begin{tikzpicture}[node distance=8mm, every node/.style={font=\sffamily\small}]
  \node[draw, rounded corners, fill=lightblue, minimum width=40mm, minimum height=8mm] (ops) {運用は「動かし続ける工学」である};
  \node[draw, rounded corners, below left=of ops, minimum width=24mm] (plan) {計画};
  \node[draw, rounded corners, below=of ops, minimum width=24mm] (deliver) {提供};
  \node[draw, rounded corners, below right=of ops, minimum width=24mm] (control) {制御};
  \draw[-{Latex}] (ops) -- (plan);
  \draw[-{Latex}] (ops) -- (deliver);
  \draw[-{Latex}] (ops) -- (control);
  \end{tikzpicture}
  \end{center}
\end{titlepage}

\tableofcontents
\newpage

\section*{全体概要：この章で学ぶこと}
\addcontentsline{toc}{section}{全体概要：この章で学ぶこと}

この章は、ソフトウェアを本番環境で安定して動かし続けるための知識領域を扱う。運用は、開発が終わった後に突然始まる作業ではない。新しいソフトウェアを作る意思決定が行われた時点から、運用計画、保守計画、サービス水準、容量、障害時の復旧、セキュリティ、監視、自動化を考え始める必要がある。

\begin{pointbox}{到達目標}
この章を読み終えると、次のことができるようになる。
\begin{itemize}
  \item ソフトウェアエンジニアリング運用の範囲を説明できる。
  \item 運用計画、運用提供、運用制御の三つのプロセス群を区別できる。
  \item 配備とリリース、インシデントと問題、バックアップとフェイルオーバーの違いを説明できる。
  \item 可用性、継続性、サービス水準、容量、リスク管理の関係を説明できる。
  \item DevOps、IaC、SRE、継続的デリバリー、自動テスト、監視とテレメトリが運用にどう関係するかを説明できる。
\end{itemize}
\end{pointbox}

\todoprompt{06章の全体地図}{中央に「ソフトウェアエンジニアリング運用」を置き、周囲に「基礎」「計画」「提供」「制御」「実務上の考慮」「ツール」を配置する。各領域から主要語として「配備」「監視」「変更」「障害」「自動化」「サービス水準」を枝で示す。}

\section*{最初に必要な用語}
\addcontentsline{toc}{section}{最初に必要な用語}

\begin{longtable}{L{35mm}L{115mm}}
\toprule
用語 & 初学者向けの説明 \\
\midrule
ソフトウェアエンジニアリング運用 & ソフトウェアを配備し、動かし、監視し、障害や変更に対応し、利用者へ安定したサービスを提供し続ける活動である。 \\
運用担当技術者 & 運用サービス、配備、監視、障害対応、容量管理、データ保護などを実行または自動化する技術者である。 \\
開発運用連携（DevOps） & 開発、保守、運用の間の壁を低くし、共通のプロセスとツールで素早く安全にソフトウェアを改善する考え方である。 \\
インフラストラクチャのコード化（IaC） & サーバ、ネットワーク、設定、環境構築などを手作業ではなくコードとして表し、再現可能に管理する方法である。 \\
プラットフォーム工学 & 開発者が自分で環境構築、配備、監視などを使えるように、内部向けの共通基盤を作る活動である。 \\
サイト信頼性エンジニアリング（SRE） & 可用性、性能、遅延、容量、緊急対応などを、ソフトウェア工学の方法で改善する運用の考え方である。 \\
サービス水準合意（SLA） & サービスの可用性、応答時間、復旧時間などについて、提供者と利用者の間で合意した水準である。 \\
インシデント & サービスの正常な運用を妨げる出来事である。障害、性能劣化、誤設定、セキュリティ事象などが含まれる。 \\
問題 & 一つまたは複数のインシデントの背後にある原因、または原因候補である。 \\
テレメトリ & ソフトウェアやインフラから収集する運用データである。ログ、実行トレース、CPU使用量、メモリ使用量、利用者操作などが含まれる。 \\
\bottomrule
\end{longtable}

\section*{導入：運用が重要な理由}
\addcontentsline{toc}{section}{導入：運用が重要な理由}

ソフトウェアエンジニアリング運用とは、ソフトウェアまたはシステムの完全性と安定性を保ちながら、配備し、運用し、支援するために必要な活動と作業である。ここでいう運用は、対象環境へのソフトウェアの配備と設定、利用中の監視と管理、欠陥対応、環境変更への対応、新しい利用者要求への対応を含む。

従来、運用部門や計算センターは、開発部門から分離された組織として扱われることが多かった。この分離は、引き継ぎの遅れ、手作業による設定ミス、責任境界の曖昧化、リリースの遅延を生みやすい。現在は、開発、保守、運用を近づけ、DevOps、IaC、プラットフォーム工学、SREなどを使って、変更を早く、かつ安全に本番へ届ける考え方が広がっている。

この章で重要な視点は二つである。第一に、運用は「最後に引き受ける部門作業」ではなく、開発初期から設計されるサービスである。第二に、現代の運用では、環境、設定、配備、監視、試験、復旧をできるだけコードと自動化で表すことで、再現性、標準化、セキュリティ、変更管理、拡張性を高める。

\begin{pointbox}{重要}
運用の品質は、利用者が直接体験するソフトウェアの品質である。コードが正しくても、配備が失敗し、監視できず、復旧できず、容量が足りなければ、利用者から見たサービスは失敗である。
\end{pointbox}

\section{ソフトウェアエンジニアリング運用の基礎}

\subsection{ソフトウェアエンジニアリング運用の定義}

ソフトウェアエンジニアリング運用とは、ソフトウェア製品、情報技術インフラ、システムソフトウェア、アプリケーションソフトウェアが、開発中、保守中、実運用中に適切に動作するように使われる知識、技能、プロセス、ツールの総体である。

運用担当技術者は、単にサーバを見張る人ではない。コンテナや仮想サーバの準備、配備、設定、支援を行う。必要な環境をオンデマンドで提供し、継続的インテグレーション、試験、配備、監視をサービスとして提供する。障害を診断し、問題と解決策を記録し、優先順位を付け、影響を評価する。さらに、セキュリティ、データ保護、フェイルオーバー、容量、ストレージ、データベース性能、技術仕様書の整備にも関わる。

\begin{longtable}{L{45mm}L{105mm}}
\toprule
役割 & 主な責任 \\
\midrule
運用担当技術者 & 運用サービスを設計し、配備、監視、容量管理、障害対応、セキュリティ、復旧を実行または自動化する。 \\
ソフトウェア技術者 & 運用担当技術者が提供する運用サービスを使い、自分のソフトウェアを配備、監視、改善する。 \\
プラットフォーム技術者 & 開発者が自律的に使える共通基盤、セルフサービス環境、配備基盤、監視基盤を作る。 \\
SRE技術者 & 可用性、性能、遅延、容量、緊急対応、変更管理を、測定と自動化に基づいて改善する。 \\
\bottomrule
\end{longtable}

\subsection{ソフトウェアエンジニアリング運用プロセス}

運用プロセスは、既存のソフトウェアシステムまたは製品を配備、設定、運用、支援し、その完全性を保つための活動である。国際標準では、運用の準備、運用の実行、運用結果の管理、利用者支援が重要な活動として扱われる。

SWEBOK 06章では、運用活動を学習しやすくするため、次の三つのプロセス群として整理できる。

\begin{longtable}{L{35mm}L{115mm}}
\toprule
プロセス群 & 含まれる活動 \\
\midrule
運用計画 & 運用計画、供給者管理、開発環境と運用環境、可用性、継続性、サービス水準、容量、バックアップ、災害復旧、フェイルオーバー、安全性とセキュリティを扱う。 \\
運用提供 & 運用テスト、検証、受け入れ、配備、リリース、ロールバック、データ移行、問題解決を扱う。 \\
運用制御 & インシデント管理、変更管理、監視、測定、追跡、レビュー、運用支援、サービス報告を扱う。 \\
\bottomrule
\end{longtable}

\todoprompt{運用プロセスの三層図}{上から「運用計画」「運用提供」「運用制御」の三つの箱を並べ、それぞれに主要活動を配置する。計画は「SLA、容量、バックアップ、セキュリティ」、提供は「テスト、配備、リリース、ロールバック」、制御は「インシデント、変更、監視、報告」とする。}

\subsection{ソフトウェア導入}

ソフトウェア導入とは、ソフトウェアまたは更新版を利用者に使える状態にするため、対象環境に配置し、設定し、確認する作業である。導入では、旧版の削除、対象環境に合わせた設定、必要なディレクトリ、登録情報、環境変数の作成などが行われる。多くの場合、これらはスクリプトで実行される。

組込みシステムでは、電子的な配布だけでなく、物理媒体を使うこともある。いずれの場合でも、導入後には検証が必要である。導入が成功したか、必要なファイルが配置されたか、設定が正しいか、最低限の動作確認ができるかを確認する。

\subsection{スクリプト化と自動化}

スクリプト化とは、繰り返し行う作業を簡易なプログラムとして記述することである。自動化とは、人が手作業で行っていた処理を、ツールやスクリプトで再現可能に実行することである。運用では、環境構築、設定変更、配備、監視、通知、ログ収集、バックアップ、復旧試験など、多くの活動が自動化の対象になる。

自動化には三つの効果がある。第一に、遅延を減らせる。第二に、手作業のばらつきや設定ミスを減らせる。第三に、障害時の反応を早め、停止時間を短くできる。さらに、自動化された手順は標準化され、運用サービスとして開発者へ提供しやすくなる。

\begin{pointbox}{実務で見るべきこと}
自動化は「便利だから行う」のではない。変更を安全に早く届けるため、同じ環境を何度でも作れるようにするため、障害時に人間の記憶へ依存しないために行うのである。
\end{pointbox}

\subsection{効果的なテストとトラブルシューティング}

運用はシステムの安定性に責任を持つ。そのため、ソフトウェアは本番配備前に十分に試験されなければならない。手作業の試験は時間がかかり、誤りや漏れが発生しやすく、大規模化しにくい。したがって、単体試験、統合試験、システム試験、受け入れ試験、回帰試験を可能な限り自動化することが重要である。

本番でエラーが見つかった場合、運用担当技術者とソフトウェア技術者は、診断を実行し、問題と解決策を記録し、優先順位を付け、影響を評価する。報告された問題が実際に再現するかを確認することも重要である。大規模ソフトウェアに対して全面的な再試験を毎回行うことは高価であるため、回帰試験や試験網羅の戦略が必要になる。

本番環境での試験は難しい。停止できない重要システムでは、カナリア試験や暗黙リリースなどを使い、一部の利用者や裏側の経路だけで新機能を観察する方法が使われる。

\subsection{性能、信頼性、負荷分散}

性能、信頼性、負荷分散は、開発初期から計画されるべきである。性能とは、応答時間や処理量が期待水準を満たすことである。信頼性とは、故障しにくく、期待どおりに動き続けることである。負荷分散とは、複数のサーバや資源に処理を振り分け、過負荷や単一点障害を避けることである。

現代の運用では、需要に応じてインフラを動的に増減させることが重視される。DevOpsの考え方を使うと、運用担当技術者は開発段階から必要なインフラサービスを用意し、ソフトウェア技術者は開発中からそれを使って試験できる。

\section{ソフトウェアエンジニアリング運用計画}

運用計画は、ソフトウェアを運用し続けるための準備を体系化する活動である。運用は多くの場合、開発期間より長く続く。したがって、運用計画では、短期の配備だけでなく、長期の保守、支援、容量、サービス水準、コスト、人員、供給者、セキュリティを考える必要がある。

運用担当技術者は、作業手順やツールを、目的に合う形で文書化する。証拠として残すべきものには、方針と計画、サービス文書、手順、プロセス、プロセス制御記録がある。

\subsection{運用計画と供給者管理}

\subsubsection{運用計画}

運用計画は、プロジェクト要求、開発者や保守担当者のニーズを、運用サービスへ翻訳する活動である。開発は数か月から数年で終わることが多いが、運用は何年も続く。そのため、資源見積もり、人員計画、教育、予算、容量、支援体制を早くから考えることが重要である。

新しいソフトウェアを開発すると決めた時点で、運用と保守の要求を考慮すべきである。概念文書、運用保守計画、運用構想（Concept of Operations, CONOPS）を作り、利用者が変更要求や問題報告をどのように出すか、誰が受け付けるか、どのように優先順位を付けるか、どのようにリリースへ反映するかを明らかにする。

\begin{longtable}{L{45mm}L{105mm}}
\toprule
運用計画で扱う項目 & 説明 \\
\midrule
役割と責任 & 新しいサービスまたは変更されたサービスを、誰が実装、運用、保守するかを決める。 \\
顧客と供給者の活動 & 顧客側と供給者側が行う作業、契約、合意を明確にする。 \\
要員と教育 & 必要な人員、技能、訓練、技術支援を計画する。 \\
プロセスとツール & 使う手順、測定、方法、ツールを定義する。 \\
容量と財務 & 将来の利用量、予算、時程、費用を見積もる。 \\
サービス受け入れ基準 & 運用開始できると判断する条件を決める。 \\
期待成果 & 運用した結果として何を測定し、どの水準を満たすかを数値で表す。 \\
\bottomrule
\end{longtable}

個々の問題報告や変更要求のレベルでは、対象となるリリース日、次版の内容、潜在的な衝突、リスク、ロールバック計画、データ移行計画、関係者への通知を扱う。

\todoprompt{運用計画とCONOPSの関係}{左に「開発開始の意思決定」、中央に「概念文書」「運用保守計画」「運用構想 CONOPS」、右に「変更要求」「問題報告」「リリース計画」「運用支援」を置き、運用計画が開発初期から始まることを矢印で示す。}

\subsubsection{供給者管理}

供給者管理とは、外部供給者の製品やサービスを、品質の高い運用サービスに結びつくように管理する活動である。ここでいう供給者には、外部委託先、クラウドサービス、インフラ提供者、プラットフォーム提供者、監視ツール提供者などが含まれる。

運用担当技術者は、供給者との関係を、製品やサービスの性質に合わせて決める。たとえば、サービスとしてのインフラ（IaaS）やサービスとしてのプラットフォーム（PaaS）を使う場合、契約、責任分界点、可用性、障害通知、データ保護、監査、終了時の移行を確認しなければならない。

\subsection{開発環境と運用環境}

ソフトウェアプロセスでは、開発環境、試験環境または品質保証環境、事前本番環境、本番環境が使われる。品質を作り込み、本番リリースのリスクを下げるには、これらの環境が一貫しており、本番環境と同期している必要がある。

\begin{longtable}{L{35mm}L{115mm}}
\toprule
環境 & 目的 \\
\midrule
開発環境 & 開発者がコードを書き、単体試験や小規模な確認を行う環境である。 \\
試験・品質保証環境 & 統合試験、品質保証、回帰試験、自動試験を実行する環境である。 \\
事前本番環境 & 本番に近い条件で、配備手順、性能、運用手順を確認する環境である。 \\
本番環境 & 実際の利用者にサービスを提供する環境である。 \\
\bottomrule
\end{longtable}

DevOpsでは、これらの環境を単一のコードリポジトリから自動的に作成することが推奨される。すべての環境が同じコード源から作られれば、環境差分による不具合を減らせる。これがインフラストラクチャのコード化につながる。

\todoprompt{環境の流れ}{左から「開発環境」「試験・品質保証環境」「事前本番環境」「本番環境」を並べ、同じコードリポジトリから各環境が自動生成される構図を描く。上部に「単一の真実の源泉」、下部に「環境差分を減らす」と注記する。}

\subsection{ソフトウェア可用性、継続性、サービス水準}

可用性とは、必要なときにサービスが利用可能である度合いである。継続性とは、大きな障害や災害が発生してもサービスを維持または復旧できる能力である。サービス水準とは、可用性、性能、応答時間、復旧時間などについて合意された水準である。

運用担当技術者は、プロジェクト初期に非機能要求として定義された可用性と継続性を満たすため、適切なインフラを計画、設計、実装、試験する。可用性は測定され、記録され、予定外の停止は調査される。サービス報告では、合意されたサービス水準に対して、実際の運用がどうだったかを示す。

SLAは、運用サービスの義務を明確にするために有効である。SLAがなければ、利用者は「常に速く、絶対に止まらない」ことを期待し、提供者は「通常は動いている」と考えるかもしれない。合意された水準がなければ、運用の成否を客観的に判断しにくい。

\subsection{ソフトウェア容量管理}

容量管理とは、現在と将来の需要に対して、十分な処理能力、記憶容量、ネットワーク、データベース性能などを確保する活動である。利用者数、トランザクション数、データ量、業務量の変化を予測し、それを具体的な要求に変換して文書化する。

容量管理は、性能と容量の問題の中心となる。新規サービスや変更サービスを設計するとき、想定負荷を見積もり、必要な資源をモデル化する。容量計画には、現在の性能、将来要求、費用を伴う選択肢、SLA達成のための推奨策を含める。

\begin{pointbox}{例}
動画配信サービスで、通常時は一日10万人が利用するが、イベント配信時には同時接続が10倍になると予測される場合、容量管理ではピーク時のサーバ数、ネットワーク帯域、キャッシュ戦略、負荷分散、監視しきい値を計画する必要がある。
\end{pointbox}

\subsection{バックアップ、災害復旧、フェイルオーバー}

バックアップとは、データ、文書、ソフトウェアを復旧できるように複製して保管することである。災害復旧とは、大規模障害や災害後にサービスを復元する活動である。フェイルオーバーとは、主系が故障したときに、予備系へ切り替えてサービスを継続する仕組みである。

復旧を成功させるには、どのバックアップ方式を使うか、どの頻度で復元点を作るか、どこに保管するか、どれくらい保持するかを決める必要がある。完全バックアップは全体を保存する方式であり、増分バックアップは前回からの差分を保存する方式である。

重要なのは、バックアップを取ることではなく、復旧できることである。災害復旧とフェイルオーバーは、計画され、定期的にリハーサルされる必要がある。ソフトウェア側にも障害処理の論理が必要であり、これは開発時に計画されなければならない。

\todoprompt{バックアップ・災害復旧・フェイルオーバーの流れ}{通常運用、バックアップ作成、障害発生、フェイルオーバー、復旧、通常運用へ戻る、という時系列を描く。各段階で「復元点」「予備系」「切替」「検証」を示す。}

\subsection{ソフトウェアとデータの安全性、セキュリティ、完全性、保護、統制}

運用では、情報セキュリティをすべてのサービス活動の中で管理する必要がある。これは、情報の安全性と可用性に関するリスク評価を行い、統制を実施することである。

運用で考えるべき統制には、経営層による情報セキュリティ方針の定義、役割と責任の割当、方針の有効性を監視する責任者、セキュリティ教育、全職員への方針周知、リスク評価と統制実装に関する専門的支援、変更が統制を壊さないことの確認、セキュリティインシデントの報告と対応が含まれる。

DevOpsの発展に伴い、開発運用セキュリティ連携（DevSecOps）が重要になっている。DevSecOpsでは、セキュリティを後工程で追加するのではなく、開発と運用の最初から組み込み、セキュリティ問題の検出と修正をできるだけ早く自動化する。

\todoprompt{DevOps・IaC・SRE・プラットフォーム工学の関係}{中央に「安定した継続的な運用」を置き、周囲に「DevOps」「IaC」「SRE」「プラットフォーム工学」「DevSecOps」を配置する。それぞれの矢印に「協働」「再現性」「信頼性」「セルフサービス」「早期セキュリティ」と注記する。}

\section{ソフトウェアエンジニアリング運用提供}

運用提供は、計画された運用サービスを実際に実行し、ソフトウェアを本番利用に耐える形で届ける活動である。ここでは、運用テスト、配備、リリース、ロールバック、データ移行、問題解決が中心となる。

\subsection{運用テスト、検証、受け入れ}

運用テストは、ソフトウェアが本番環境または本番に近い環境で運用可能かを確認する活動である。検証は「仕様どおりか」を確認する活動であり、受け入れは「利用者や運用者が使えると判断できるか」を確認する活動である。

現代の運用では、テストは最後に一回だけ行うものではない。テスト駆動開発（TDD）や受け入れテスト駆動開発（ATDD）を使い、開発中から継続的に検証を行う。DevOpsでは、試験サービスや試験ツールを自動化し、開発者が変更のたびに素早く結果を得られるようにする。

\subsection{配備・リリース工学}

配備とは、指定された版のソフトウェアを、特定の環境へインストールすることである。リリースとは、機能または機能集合を、顧客や利用者が実際に使える状態にすることである。配備とリリースは似ているが、同じではない。配備は技術的な配置であり、リリースは利用可能にする判断である。

\begin{longtable}{L{35mm}L{115mm}}
\toprule
概念 & 説明 \\
\midrule
配備 & コード、設定、データベース、サーバ、コンテナなどを対象環境へ置き、動作可能にする技術的活動である。 \\
リリース & 配備済みの機能を、すべての利用者または一部の利用者に使えるようにする活動である。ビジネス判断を含む。 \\
\bottomrule
\end{longtable}

従来は、開発チームが完成物を運用チームに引き渡し、運用チームが手作業で配備することが多かった。この方法は時間がかかり、品質面でもリスクがある。DevOpsでは、コードの包装、設定ファイル生成、サーバ再起動、サーバとデータベースの設定、ソフトウェアのインストール、アプリケーション起動、スモークテストを自動化する。

リリース戦略には、環境に基づく戦略とアプリケーションに基づく戦略がある。環境に基づく戦略では、新版を事前本番環境や別環境へ配備してから切り替える。アプリケーションに基づく戦略では、機能トグルのような設定により、コード中の特定機能を有効化または無効化する。

カナリアリリースは、変更を一部の範囲にだけ短期間配備し、その結果を評価してから全面配備を判断する方法である。新しいソフトウェアが異常を示した場合、自動的にロールバックできる構成も重要である。

\todoprompt{配備とリリースの違い}{左に「配備：ソフトウェアを環境へ置く」、右に「リリース：利用者に使わせる」を並べる。中央に「本番環境に配備済みだが未リリースの機能」という中間状態を描く。機能トグルとカナリアリリースを注記する。}

\subsection{ロールバックとデータ移行}

ロールバックとは、新しい版のソフトウェアやデータベースが問題を起こしたとき、正しく動く以前の状態へ戻す活動である。データ移行とは、ソフトウェア変更に合わせてデータ構造や内容を移す活動である。

本番配備前には、ロールバック手順が計画され、練習されるべきである。新しい版が障害や性能劣化を起こした場合、素早く以前の状態へ戻せる必要がある。DevOpsでは、この過程を自動化し、監視結果に応じて自動的に戻すことも可能である。

データ移行がある場合、単にプログラムを戻すだけでは不十分である。データの形式が変わっていると、旧版が新しいデータを読めないことがある。したがって、ロールバック計画には、データをどの状態に戻すか、どこまでの処理を再実行するか、利用者にどの影響があるかを含める必要がある。

\todoprompt{ロールバックとデータ移行の時系列}{横軸を時間とし、「旧版稼働」「新版配備」「データ移行」「監視」「異常検知」「ロールバック」「データ復元または補正」を並べる。ソフトウェア状態とデータ状態の二段で示す。}

\subsection{問題解決}

問題解決は、ソフトウェアやシステムのインシデントと問題の原因を識別し、分析し、事業への混乱を最小化する活動である。繰り返し発生する本番問題には、ソフトウェア、インフラ、設定、データベース、ネットワーク、人間の手順など、複数の原因が絡むことがある。そのため、問題解決には、ソフトウェア技術者、運用担当技術者、セキュリティ担当、データベース担当、製品責任者などの多職種チームが必要になる場合がある。

問題解決では、監視、ログ、プロファイル、実行トレース、設定履歴を使って、何が起きたか、いつ起きたか、どの変更と関係するかを調べる。問題の根本原因がわかれば、再発防止策を実装する。

\section{ソフトウェアエンジニアリング運用制御}

運用制御は、運用中のサービスを観察し、制御し、問題に対応し、変更を管理し、結果を報告する活動である。運用は一度設定すれば終わりではない。利用者、負荷、脅威、法規制、事業要求、技術基盤は変化する。その変化を制御しなければ、サービス品質は低下する。

\subsection{インシデント管理}

インシデント管理とは、ソフトウェアインシデントを記録し、優先順位を付け、事業影響を評価し、解決、エスカレーション、終了まで管理するプロセスである。

現代のDevOpsでは、アラートとログを用いた監視を自動化し、小さなインシデントが大きな障害になる前に検出する。インシデントが発生したら、原因を分析し、事後検討を行い、同じインシデントを繰り返さないための対策を実施する。

\subsection{変更管理}

変更管理とは、すべての変更を評価、承認、実装、レビューされた状態で行うためのプロセスである。変更要求は記録され、緊急、至急、重大、軽微などに分類される。変更管理では、変更のリスク、ロールバック戦略の必要性、関係者への通知、変更スケジュールを扱う。

従来型の開発では、多くの変更をまとめて定期リリースとして提供することが多かった。DevOpsでは、小さな変更単位を独立に、必要に応じて届けることを目指す。そのためには、ソフトウェアが小さく独立した配備に向くように設計されている必要がある。

\subsection{監視、測定、追跡、レビュー}

運用では、容量、継続性、可用性を監視する。DevOpsの考え方では、希望的観測に頼ってはならない。システム品質と運用状態を、証拠に基づいて判断する必要がある。

\begin{longtable}{L{55mm}L{95mm}}
\toprule
測定対象 & 意味 \\
\midrule
本番システム監視と製品テレメトリ & 本番で何が起きているかを示す。ログ、トレース、利用者行動、資源使用量などを含む。 \\
リリース前後の検証結果 & リリースが期待どおりに動くか、リリース後に品質が悪化していないかを示す。 \\
利用者活動と資源使用 & どの機能が使われ、どの資源が消費されているかを示す。 \\
影響分析結果 & ある変更がどの機能、構成、利用者に影響するかを示す。 \\
依存関係 & サービス運用に必要な内部依存と外部依存を示す。 \\
構成変更 & 承認済み配備以外の設定変更や環境差分を検出する。 \\
セキュリティと耐障害性 & 攻撃、障害、復旧に対する能力を示す。 \\
\bottomrule
\end{longtable}

\todoprompt{監視とテレメトリのデータ層}{下から「インフラ層」「オペレーティングシステム層」「アプリケーション層」「利用者行動層」を積み上げ、各層からログ、トレース、メトリクス、イベントが収集され、分析基盤、ダッシュボード、アラートへ流れる図にする。}

\subsection{運用支援}

運用支援は、ソフトウェア製品を意図された環境で動かし、顧客や利用者を支援する活動である。運用支援はプロジェクトの計画段階から始まり、実行段階では監視、イベント対応、インシデント対応、問い合わせ対応、手順実行、運用文書更新を行う。

運用支援活動はSLAに記述されることが多い。たとえば、サポート時間、初動応答時間、障害の重大度分類、復旧目標、報告形式、エスカレーション経路などである。

\subsection{サービス報告}

サービス報告は、意思決定に必要な、合意済みで、適時で、信頼でき、正確な情報を作る活動である。報告は、運用サービスがどのように機能したか、サービス水準を満たしたか、どの問題が発生したか、今後の傾向がどうかを示す。

典型的な報告項目には、サービス水準目標に対する性能、セキュリティ侵害、トランザクション量、資源使用量、インシデントと失敗、傾向情報、満足度分析がある。測定を自動化し、収集、分析、報告、制御を調整できる枠組みとツールを選定することが重要である。

\todoprompt{インシデント・問題・変更管理の関係}{「インシデント発生」から「記録」「一次対応」「問題分析」「根本原因」「変更要求」「変更承認」「修正配備」「再発防止」へつながる流れを描く。インシデント管理、問題解決、変更管理の担当範囲を色ではなく線種や枠で区別する。}

\section{実務上の考慮}

\subsection{インシデントと問題の予防}

インシデントや問題を予防するには、運用プロセス全体を可能な限り自動化し、試験を継続的に統合する必要がある。さらに、製品テレメトリを適切な分析技法と組み合わせ、問題を早期に検出する。

データは、アプリケーション層、オペレーティングシステム層、インフラ層など、製品スタックのすべての層から収集されるべきである。テレメトリは、問題の兆候を発見するだけでなく、問題源を特定するための基盤になる。

\subsection{運用リスク管理}

運用担当技術者は、可用性、拡張性、セキュリティ、個人情報漏えい、誤配備、容量不足、外部サービス停止など、多くのリスクを管理する。継続的リスク管理では、運用状態を常に監視し、リスクを検出したら自動的にアラートを出す。

何をアラートにするかは、製品責任者やソフトウェア技術者と相談して、リスク許容度を定める必要がある。すべての小さな変化で警報を出すと、警報疲れが起こる。逆に、重大な兆候を見逃すと、障害や事故につながる。

\subsection{ソフトウェアエンジニアリング運用の自動化}

現代の運用では、自動化が中心的役割を持つ。環境構築、ユーザアカウント作成、サーバ準備、実行時設定変更、配備、試験、デバッグ支援などを自動化できる。ソフトウェア技術者は、アプリケーション自動化と運用自動化を結びつけることで大きな効果を得る。

運用自動化の傾向は、複雑さを減らし、インフラ準備を速め、開発者へ運用サービススクリプトを提供し、アプリケーション定義、配備、試験作業を自動化する方向にある。

\todoprompt{CI/CDと継続的テストのパイプライン}{左から「コード変更」「ビルド」「単体試験」「統合試験」「パッケージ」「事前本番配備」「受け入れ試験」「本番配備」「監視」「フィードバック」を並べる。各段階に自動化の歯車アイコンを置き、監視結果がコード変更へ戻るループを描く。}

\subsection{小規模組織における運用}

小規模組織、とくに25人以下の非常に小さい組織では、大規模組織向けの標準をそのまま適用することが難しい。標準が要求する文書、役割、手順、監査が過大になることがある。この場合、小規模組織向けに調整されたISO/IEC 29110系列のような指針が有用である。

重要なのは、小規模であることを理由に運用を無視しないことである。むしろ、少人数だからこそ、自動化、標準化、簡潔な手順、明確な責任分担が重要になる。小規模組織では、最小限のSLA、バックアップ手順、インシデント対応、変更管理、監視ダッシュボードから始めるとよい。

\section{ソフトウェアエンジニアリング運用ツール}

運用ツールは、人的資源を効率よく使い、品質と応答速度を高めるために重要である。適切に実装された自動化は、手作業より速く、容易で、信頼性が高い。DevOpsでは、開発、保守、運用の資源と手順を結びつけ、継続的インテグレーション、継続的デリバリー、継続的テスト、継続的デプロイを支援する。

継続的デリバリーは、新しいシステムやソフトウェアを頻繁に試験環境や事前本番環境へ届ける実践である。継続的テストは、ソフトウェアライフサイクルの各段階で品質を評価する実践である。継続的デプロイは、自動検証を通過した変更を本番へ自動配備するプロセスである。

\subsection{コンテナと仮想化}

コンテナと仮想化は、アプリケーションの拡張性を高め、複数のサーバやクラウド提供者にまたがる配備を標準化するために使われる。コンテナは、アプリケーションと実行に必要な依存関係をまとめ、環境差分を減らす。仮想化は、物理サーバ上に複数の仮想サーバを作る。

運用担当技術者は、プロジェクトの規模、複雑さ、柔軟性、セキュリティ、監視要求に基づいてツールを選ぶ。コンテナ管理ツール、すなわちオーケストレータは、多数のコンテナの配備、再起動、スケーリング、ネットワーク、監視を支援する。

\subsection{配備}

配備ツールは、さまざまな環境へのソフトウェア配備を管理する。配備記述ファイル、設定ファイル、環境変数、シークレット管理、コンテナイメージ、クラウド資源、データベース変更などを扱う。

実務では、一つのツールだけで全体をまかなうことは少ない。ビルド、パッケージ、インフラ作成、設定管理、データベース移行、監視、アラート、ロールバックなど、複数のツールを組み合わせて配備パイプラインを構築する。

\subsection{自動テスト}

開発者へ速く継続的なフィードバックを提供するため、テストはソフトウェア提供プロセス全体で可能な限り自動化されるべきである。単体試験、統合試験、システム試験、利用者受け入れ試験を含む試験戦略を定義し、それぞれを支援するツールを選ぶ。

自動テストは、単に欠陥を検出するだけではない。変更が安全に本番へ進めるかを判断する門番である。自動テストが不十分だと、配備を速くするほど障害を速く広げてしまう。

\subsection{監視とテレメトリ}

監視とテレメトリは、運用の中核である。監視は、サービスが正常かどうかを観察する活動である。テレメトリは、システムからデータを収集して分析に使う仕組みである。

監視では、アプリケーション層のログ、オペレーティングシステム層の実行トレース、サーバ層のCPU使用量やメモリ使用量などを収集する。収集されたデータには、統計分析や機械学習を使って意味を抽出できる。最終的には、ダッシュボードにより、開発者、運用者、製品責任者、経営層など、異なる利害関係者に必要な情報を見せる。

\todoprompt{運用ツールの分類図}{中央に「運用ツール」を置き、周囲に「コンテナ・仮想化」「配備」「自動テスト」「監視・テレメトリ」「設定管理」「セキュリティ」を配置する。各カテゴリに代表的な支援内容を短く添える。}

\section{トピックと参考文献の対応}

この表は、SWEBOK 06章のトピックと主要参考資料を、学習用に簡略化したものである。

\begin{longtable}{L{70mm}L{80mm}}
\toprule
トピック & 主な参考資料 \\
\midrule
ソフトウェアエンジニアリング運用の基礎 & ISO/IEC/IEEE 20000-1、ISO/IEC/IEEE 12207、The DevOps Handbook \\
運用プロセス & ISO/IEC/IEEE 20000-1、ISO/IEC/IEEE 12207、ISO/IEC/IEEE 32675 \\
ソフトウェア導入、スクリプト化、自動化 & The DevOps Handbook、ISO/IEC/IEEE 20000-1 \\
テストとトラブルシューティング & The DevOps Handbook、ソフトウェアテスト知識領域 \\
性能、信頼性、負荷分散 & ISO/IEC/IEEE 20000-1、SRE文献 \\
運用計画と供給者管理 & ISO/IEC/IEEE 20000-1、ISO/IEC/IEEE 12207 \\
開発環境と運用環境 & The DevOps Handbook、IaC関連標準 \\
可用性、継続性、サービス水準 & ISO/IEC/IEEE 20000-1、ソフトウェア保守知識領域 \\
容量管理 & ISO/IEC/IEEE 20000-1 \\
バックアップ、災害復旧、フェイルオーバー & ISO/IEC/IEEE 20000-1、SRE文献 \\
安全性、セキュリティ、完全性、保護、統制 & ISO/IEC/IEEE 20000-1、ISO/IEC/IEEE 32675、ソフトウェアセキュリティ知識領域 \\
配備・リリース工学 & The DevOps Handbook、Continuous Delivery、ソフトウェア構成管理知識領域 \\
ロールバックとデータ移行 & The DevOps Handbook、ISO/IEC/IEEE 12207 \\
インシデント管理、変更管理、問題解決 & ISO/IEC/IEEE 20000-1、The DevOps Handbook \\
監視、測定、追跡、レビュー & The DevOps Handbook、The Art of Monitoring \\
小規模組織における運用 & ISO/IEC 29110系列 \\
運用ツール & The DevOps Handbook、The Art of Monitoring、ISO/IEC/IEEE 32675 \\
\bottomrule
\end{longtable}

\section{参考文献}

\begin{enumerate}[label={\arabic*.}]
  \item ISO/IEC/IEEE 20000-1:2013, Information technology -- Service management -- Part 1: Service management systems requirements.
  \item G. Kim, J. Humble, J. Debois, J. Willis, and N. Forsgren, \textit{The DevOps Handbook: How to Create World-Class Agility, Reliability and Security in Technology Organizations}, 2nd ed., IT Revolution Press, 2021.
  \item ISO/IEC/IEEE 12207:2017, Systems and software engineering -- Software Life Cycle Processes.
  \item ISO/IEC/IEEE 32675:2022, Information Technology -- DevOps: Building Reliable and Secure Systems Including Application Build, Package and Deployment.
  \item ISO/IEC/IEEE 24765:2017, Systems and Software Engineering -- Vocabulary, 2nd ed., 2017.
  \item B. Beyer, C. Jones, J. Petoff, and N. R. Murphy, \textit{Site Reliability Engineering: How Google Runs Production Systems}, O'Reilly Media, 2016.
  \item ISO/IEC CD 29110-5-5:2023, Systems and software engineering -- Lifecycle profiles for Very Small Entities, Part 5-5: Agile/DevOps guidelines.
  \item J. Humble and D. Farley, \textit{Continuous Delivery: Reliable Software Releases through Build, Test, and Deployment Automation}, Pearson Education, 2010.
  \item J. Turnbull, \textit{The Art of Monitoring}, James Turnbull, 2016.
\end{enumerate}

\section{画像 TODO 一覧}

本文に配置した画像TODOを再掲する。実際に図を作る場合は、各TODOのプロンプトをそのまま画像生成に使うと、A4資料に合う図を作りやすい。

\begin{enumerate}
  \item 06章の全体地図。
  \item 運用プロセスの三層図。
  \item 運用計画とCONOPSの関係。
  \item 環境の流れ。
  \item バックアップ・災害復旧・フェイルオーバーの流れ。
  \item DevOps・IaC・SRE・プラットフォーム工学の関係。
  \item 配備とリリースの違い。
  \item ロールバックとデータ移行の時系列。
  \item 監視とテレメトリのデータ層。
  \item インシデント・問題・変更管理の関係。
  \item CI/CDと継続的テストのパイプライン。
  \item 運用ツールの分類図。
\end{enumerate}

\section{章末確認問題}

\begin{enumerate}
  \item ソフトウェアエンジニアリング運用は、単なる本番配備とどう違うか。
  \item 運用計画、運用提供、運用制御の違いを説明せよ。
  \item 運用担当技術者とソフトウェア技術者の役割の違いを説明せよ。
  \item IaCが環境差分を減らす理由を説明せよ。
  \item 配備とリリースの違いを、具体例を使って説明せよ。
  \item カナリアリリースがリスクを減らす理由を説明せよ。
  \item ロールバック計画にデータ移行が含まれるべき理由を説明せよ。
  \item インシデント管理、問題解決、変更管理の関係を説明せよ。
  \item 監視とテレメトリが、インシデント予防に役立つ理由を説明せよ。
  \item 小規模組織では、どの運用活動から始めるのが現実的かを述べよ。
\end{enumerate}

\begin{pointbox}{解答の方向性}
1は、配備、設定、監視、障害対応、変更管理、容量管理、バックアップ、復旧、セキュリティを含む点を述べる。2は、計画が準備、提供が実行、制御が監視と管理である点を述べる。5は、配備は環境へ置くこと、リリースは利用者に使わせることと説明する。8は、インシデントを記録し、問題として原因を分析し、変更管理で修正を安全に入れる流れとして説明できる。
\end{pointbox}

\end{document}
